% ============================================================
% 术语对照表（Glossary）— 模板
% 中文翻译的唯一术语标准
% 全文必须统一使用，不得出现同义异译
%
% 英文术语首次出现时保留英文括注；之后全文仅使用中文。
%
% 用法：使用前按论文领域替换具体术语，
%       保留"不翻译的项"和"翻译原则"部分。
% ============================================================

% --- 核心数学对象 ---
% half-derivation                 1/2-导子
% N_n                             幂零李代数 N_n（严格上三角矩阵）

% --- 算子与空间 ---
% evaluation map                  求值映射 ev
% evaluation matrix               求值矩阵 M
% constraint operator             约束算子 \widetilde L

% --- 证明架构：约束化简层 ---
% constraint reduction            约束化简
% chain bridge                    链式桥接（首次）/ 链桥（L2）
% image restriction               像限制（L5）
% parameter count                 参数统计（L7）

% --- 导子分类 ---
% centered derivation             中心化（centered）导子（首次）
%                                 中心化导子（后续简称）
% boundary cocycle                边界 cocycle

% --- 信道变量 ---
% a-channel                       a-信道
% channel variable                信道变量

% --- Dynkin 图相关 ---
% Dynkin diagram                  Dynkin 图
% Dynkin endpoint                 Dynkin 端点

% --- 矩阵相关 ---
% evaluation matrix               求值矩阵
% rank                            秩

% --- 上下界 ---
% force / forces                 （按语境：推出 / 必有 / 可知 / 必须满足）

% --- 不翻译的项（保留原样，与英文版一致）---
% coeffOf_cond                    保留
% finrank_halfDer                 保留

% --- 通用数学术语 ---
% characteristic                  特征
% field                           域
% linear map                      线性映射
% vanish / vanishes               为零 / 消失

% --- 论文结构 ---
% Section                         节 (§)
% forward reference                前向引用

% ============================================================
% 翻译原则（翻译时必须遵守）
% ============================================================

% 0. 英文术语首次出现时保留英文括注；之后全文仅使用中文。

% 1. 英文定理名、Lean 定理名、文件名一律保留英文，不翻译。

% 2. force 在正文中不机械翻译为"迫使"。按语境选择：
%    "推出" / "必有" / "从而得到" / "必须满足" / "可知"。

% 3. chain bridge：首次出现写"链式桥接（chain bridge）"，
%    后续写"链桥"。

% 4. centered：首次出现写"中心化（centered）"，后续写"中心化"。

% 5. boundary cocycle：全文保留 cocycle，不翻译为"余循环"。

% 6. 两阶段翻译流程：
%    第一阶段：保持 LaTeX 结构、公式、标签、引用完全不变，只翻正文。
%    第二阶段：数学论文风格润色（非英语直译）。

% 7. 不逐句翻译英文，按中文数学论文的表达方式重写。
%    允许：调整句序、合并短句、被动语态改为主动语态。
%    要求：数学含义保持完全一致。
